\begin{lemma}
Let $D$ be sufficiently large, let $\mathcal F\in H(3,3,D)$, and let
$f\in E(\mathcal F)$ maximize $b_{L(\mathcal H),3D}(e)$ over all
$\mathcal H\in H(3,3,D)$ and all $e\in E(\mathcal H)$.  Then every vertex
$v\in f$ has $\deg_{\mathcal F}(v)=D$, and $f$ intersects exactly
$3(D-1)$ other edges of $\mathcal F$.
\end{lemma}
\begin{proof}
This is Lemma~\ref{lem: structure1} of the paper specialized to $k=3$ and
$t=3$.  We give the local perturbation argument and check the $3$-simple
membership directly.

First suppose a vertex $v_1\in f$ has degree less than $D$.  Add a fresh edge
$g=\{v_1,x_1,x_2\}$ using two new vertices, and keep the same distinguished
edge $f$.  By \noderef{HypergraphClass}, to verify that the perturbed
hypergraph remains in $H(3,3,D)$ we must check the three component
conditions.  It is $3$-uniform by \noderef{UniformHypergraph}, since $g$ has
three vertices and all old edges are unchanged.  It is $3$-simple by
\noderef{TSimpleHypergraph}, because old-old intersections are unchanged and
the new edge can meet an old edge only in the single old vertex $v_1$.  It has
maximum degree at most $D$ by \noderef{MaxDegreeAtMost} and
\noderef{HypergraphDegree}: the fresh vertices have degree $1$, the old
vertices other than $v_1$ keep their degrees, and the degree of $v_1$ rises
from a value strictly below $D$ to a value at most $D$.

Let $d=\deg_{L(\mathcal F)}(f)$ and let $P,T$ be the old values
$P_{L(\mathcal F)}(f),T_{L(\mathcal F)}(f)$ from
\noderef{IndependentPairCount} and \noderef{IndependentTripleCount}.  By
\noderef{LineGraphOfHypergraph}, the new edge $g$ is a new neighbor of $f$,
and no old adjacency to $f$ changes; hence the line-graph degree of $f$
becomes $d+1$.  The new independent pairs in the neighborhood of $f$ are
precisely the pairs $\{g,h\}$ with $h$ an old neighbor of $f$ disjoint from
$g$.  Since $x_1,x_2$ are fresh, this means exactly that $h$ does not contain
$v_1$.  There are $d-(\deg_{\mathcal F}(v_1)-1)$ such old neighbors, because
\noderef{HypergraphDegree} counts $\deg_{\mathcal F}(v_1)-1$ edge indices
other than $f$ containing $v_1$.  Thus
\[
P_{\rm new}=P+d-\deg_{\mathcal F}(v_1)+1 .
\]
Similarly, a new independent triple must contain $g$ and an old independent
pair of neighbors neither of whose two edges contains $v_1$.  Therefore
\[
T_{\rm new}=T+P-Q,
\]
where $Q$ is the number of old independent pairs in the neighborhood of $f$
that involve an edge containing $v_1$; here $0\le Q\le P$.

Substituting these three changes into \noderef{LocalBParameter} with
$M=3D$ gives
\[
\begin{aligned}
b_{\rm new}(f)-b_{\rm old}(f)
&=\frac1{3D}
-\frac{d-\deg_{\mathcal F}(v_1)+1}{(3D)^2}
+\frac{P-Q}{(3D)^3}  \\
&\ge \frac1{3D}-\frac{d+1}{(3D)^2}.
\end{aligned}
\]
Finally $d$ is at most the number of incidences between $f$ and edges other
than $f$.  The vertex $v_1$ contributes at most $D-2$ such incidences, because
$\deg_{\mathcal F}(v_1)<D$ and one incident edge is $f$, and each of the other
two vertices of $f$ contributes at most $D-1$.  Hence $d+1\le 3D-3<3D$, so
the displayed lower bound is positive.  This contradicts maximality of
$(\mathcal F,f)$.  Thus every vertex of $f$ has degree $D$.

Now suppose $f$ intersects fewer than $3(D-1)$ other edges.  The three
vertices of $f$ have $3(D-1)$ incidences with edges other than $f$, so some
neighboring edge $e$ contains at least two vertices of $f$, say $v_1,v_2$.
Replace $e$ by
\[
e_1=\{v_1,x_1,x_2\},\qquad
e_2=(e\cup\{x_1\})\setminus\{v_1\},
\]
where the $x_i$ are fresh, and keep $f$ as the distinguished edge.  Again
\noderef{HypergraphClass} reduces membership in $H(3,3,D)$ to the component
conditions.  Both new edges have size three, so \noderef{UniformHypergraph}
is preserved.  The replacement preserves $3$-simplicity by
\noderef{TSimpleHypergraph}: old-old intersections not involving $e$ do not
change; $e_1$ uses only one old vertex; $e_2$ is obtained from $e$ by replacing
the old vertex $v_1$ by the fresh vertex $x_1$; and
$e_1\cap e_2=\{x_1\}$.  Hence every new intersection has size at most the
corresponding old intersection size, or else is visibly a singleton.  The
maximum-degree bound from \noderef{MaxDegreeAtMost} and
\noderef{HypergraphDegree} is also preserved: $v_1$ loses $e$ and gains
$e_1$, every old vertex of $e\setminus\{v_1\}$ loses $e$ and gains $e_2$,
all other old degrees are unchanged, and the fresh vertices have degrees at
most two.

Let $d=\deg_{L(\mathcal F)}(f)$ before the replacement.  By
\noderef{LineGraphOfHypergraph}, the old neighbor $e$ of $f$ is removed and
the two new edges $e_1,e_2$ are both neighbors of $f$, so the line-graph
degree of $f$ becomes $d+1$.  We next compare the independent-pair counts
from \noderef{IndependentPairCount}.  Old independent pairs not involving
$e$ remain present unless one of their two old edges changes, and no old edge
except $e$ changes.  The only genuinely new independent pairs are therefore
pairs $\{e_1,h\}$ or $\{e_2,h\}$ with $h$ an old neighbor of $f$ different
from $e$.  Let
\[
A=\{h:\ h\text{ is an old neighbor of }f,\ h\ne e,\ h\cap e_1=\varnothing\}
\]
and
\[
B=\{h:\ h\text{ is an old neighbor of }f,\ h\ne e,\ h\cap e_2=\varnothing\}.
\]
The number of new pairs involving $e_1$ or $e_2$ is $|A|+|B|$; the pair
$\{e_1,e_2\}$ is not independent because the two edges meet at $x_1$.  A
neighbor $h$ lies in both $A$ and $B$ exactly when it is disjoint from both
new edges.  Since $x_1,x_2$ are fresh and
$e_1\cup e_2=e\cup\{x_1,x_2\}$ as far as old vertices are concerned, this is
equivalent to $h\cap e=\varnothing$.  Thus the double-counted elements
$A\cap B$ are precisely the old independent pairs $\{e,h\}$ that are lost
when $e$ is removed.  After those lost pairs are subtracted, the net increase
in $P_{L(\mathcal F)}(f)$ is at most
\[
|A\cup B|\le d-1<3D,
\]
because $A\cup B$ is a set of old neighbors of $f$ other than $e$.

The independent-triple count from \noderef{IndependentTripleCount} does not
decrease.  Indeed, an old independent triple not using $e$ remains unchanged.
If an old independent triple is $\{e,h_1,h_2\}$, then $h_1$ and $h_2$ are
disjoint from $e$ and from each other; because $e_2\setminus\{x_1\}\subseteq e$
and $x_1$ is fresh, the triple $\{e_2,h_1,h_2\}$ is an independent triple in
the new line-graph neighborhood of $f$.  This replacement is injective and its
image is disjoint from the unchanged triples not using $e$.

Using \noderef{LocalBParameter} with $M=3D$, the degree contribution to
$b_{L(\mathcal F),3D}(f)$ increases by $1/(3D)$, the pair contribution
decreases the parameter by less than $3D/(3D)^2=1/(3D)$, and the triple
contribution is nonnegative.  Hence the total change in the $b$-parameter is
strictly positive.  This contradicts maximality of $(\mathcal F,f)$.  Hence
$f$ meets exactly $3(D-1)$ other edges.
\end{proof}
